\section{Introduction}
Cirrhosis is a chronic condition that is currently the 11th leading cause of death in the US13. As liver disease progresses towards cirrhotic conditions, fibrosis in the tissue causes increased resistance to blood flow resulting in increased blood pressure in the portal venous system23. Peoples with portal hypertension, which is defined by a hepatic venous pressure gradient (HVPG) greater than 10 mm Hg, are at increased risk of developing gastroesophageal varices (GEVs)11,24. GEVs are a form of portosystemic collateral in which blood is shunted away from the liver to reduce pressure via reversed flow from the left gastric vein (LGV)13. This reversed flow causes the veins to dilate which makes then highly prone to rupture which can then lead to internal bleeding and eventually death11. Identification of at-risk GEVs before rupture is critical as GEVs can be treated, however, the primary method of diagnosis is a highly invasive, expensive, and uncomfortable procedure known as a esophagogastroduodenoscopy (EGD)25. Only 7\% of cirrhotic patients develop GEVs which results in EGDs having a large risk to benefit ratio14. As a result, there is a pressing need for the development of non-invasive methods for identifying and stratifying GEVs with high sensitivity and specificity.

4D flow MRI has been proposed as a potential alternative to EGDs as it is able to provide comprehensive hemodynamic information noninvasively14. However, the primary challenge with 4D flow MRI is the long scan times required to produce images with not only sufficient spatial and temporal resolution, but also minimal noise and velocity aliasing. Scan time during a clinical MRI session is at a premium so there is high interest in developing methods to accelerate the acquisition of 4D flow MRI, namely through the elimination of extraneous scans for dealing with velocity aliasing and utilization of unused data in respiratory phases that are normally discarded.

\section{Materials \& Methods}
Free-breathing, radial 4D flow scans, with low (15 cm/s), medium (30 cm/s), and high (60 cm/s) venc, will be acquired for 19 subjects from a mixed cohort of healthy subjects and cirrhotic patients at risk for GEVs on a clinical 3T scanner (Premier, GE Healthcare). Acquisition order of each venc will be randomized between subjects to prevent influence on the analysis. Scan times will be approximately 12 minutes long and ECG and respiratory bellows signals will be recorded for retrospective gating. Time-averaged and time-resolved images will be reconstructed using an iterative CG-SENSE algorithm with a 50\% expiration threshold42. Flow analysis will be performed in GTFlow (Gyrotools). For the flow analysis, 10 planes will be placed on the portal vein at each of the major branches: Superior Mesenteric Vein (SMV), Splenic Vein (SV) x2, Left Gastric Vein (LGV) (when present), Portal Vein (PV) x3, Left Portal Vein (LPV), Right Portal Vein (RPV), and Middle Portal Vein (MPV) (on subjects with a trifurcation). Phase unwrapping will be performed for the low and medium vencs with a dual venc unwrapping algorithm using the high venc scan written and with a single-step 4D Laplacian unwrapping algorithm26,43. To compare each of the methods, conservation of mass will be calculated at the confluence and bi/trifurcations as a measure of repeatability and correlation with the high venc will be calculated for as a measure of precision.

\section{Results}
Low, medium, and high venc 4D flow scans were acquired for 19 subjects. Scan parameters include: TE/TR = 3.2 ms/8.6 ms, flip angle = 14$^\circ$, spatial resolution = 1.25 mm isotropic, scan time = 12 min, number of radial projections = 16,800, reconstructed cardiac frames = 14. Time-averaged and time-resolved images were reconstructed using an iterative CG-SENSE algorithm with a 50\% expiration threshold. Segmentation and placement of analysis planes for calculation of flow parameters were accomplished in GTFlow (Gyrotools) using the time-averaged medium venc angiograms and then propagating planes to the time-resolved datasets for all vencs. Ten analysis planes were placed on the portal vein at each of the major branches: SMV, SV1, SV2, LGV, PV1, PV2, PV3, LPV, RPV, and MPV\@. Phase unwrapping was performed for the low and medium vencs with a dual venc unwrapping algorithm using the high venc scan written in MATLAB (Mathworks)43. 

Of the 19 subjects scanned, 7 were analyzed with the dual-venc method. The unwrapped medium venc dataset demonstrated the best agreement in conservation of mass and the best correlation with the high venc dataset while preserving a superior VNR (Figure 6). The unwrapped low venc dataset showed less correlation with the high venc dataset, possibly due to intravoxel velocity distribution

\section{Discussion}

\section{Conclusion}